\chapter{Классификация физического интертвинера стрелки--концевой сектор}
\label{ch:v8-physical-arrow-endpoint-intertwiner-classification-gate}

\section{Правильный вопрос после размерностного запрет}

Пространство одиннадцати меток стрелок нельзя отождествлять с
одиннадцатимерным физическим модулем состояний. Поэтому классифицируется полный
модуль
\begin{equation}
 \mathcal E_{\rm new}=\bigoplus_{e\in E_{\rm new}}
 \operatorname{Hom}(\mathcal H_{s(e)},\mathcal H_{t(e)}),
 \qquad \dim_{\mathbb C}\mathcal E_{\rm new}=36,
\end{equation}
и его эквивариантные отображения в модуль концевых секторов размерности $21$.

Использованы уже зафиксированные представления
$SU(3)_c\times SU(2)_L\times U(1)_Y$. Для стрелки $s\to t$ представление
равно $R_t\otimes R_s^*$. Размерность пространства интертвинеров вычислена
по совпадающим неприводимым компонентам с учётом их кратностей.

\section{Полный модуль неоднозначен}

Для одной физической ориентации получено
\begin{equation}
 \dim_{\mathbb C}\operatorname{Hom}_G(\mathcal E_{\rm new},H_{21})=10.
\end{equation}
После добавления вещественной структуры обратных сопряжённых стрелок размерность возрастает до
\begin{equation}
 \dim_{\mathbb C}\operatorname{Hom}_G(
 \mathcal E_{\rm new}\oplus J\mathcal E_{\rm new},H_{21})=14.
\end{equation}
Следовательно, полный физически эквивариантный коннектор не уникален.

Более того, все десять интертвинеров прямой ориентации лежат на пяти
Ходж-заглушённых рёбрах. На выбранной шестирёберной опоре прямая
ориентация не даёт ни одного отображения.

\section{Единственный канал партнёра вещественной структуры}

После вещественного завершения выбранной Ходж-опоры остаётся ровно один комплексный
интертвинер. Он принадлежит обратной стрелке
\begin{equation}
 Y_R\longrightarrow Q_L
\end{equation}
для выбранного ребра $Q_L--Y_R$. Слабое свёртывание выделяет компоненту
\begin{equation}
 (mathbf3,mathbf1,2/3),
\end{equation}
совпадающую с блоком концевого сектора $u_R$. Поэтому
\begin{equation}
 \dim_{\mathbb C}\operatorname{Hom}_G(
 J\mathcal E_* ,H_{21})=1.
 \label{eq:v8-unique-real-selected-intertwiner}
\end{equation}

Это первая физически типизированная одномерная зацепка. Однако её
максимальный операторный ранг равен трём: она соединяет только цветовой
триплетный блок $u_R$. Она не является изоморфизмом полных носителей и сама
по себе не устраняет центральную массовую метрику.

\begin{theorem}[Уникальный выбранный канал вещественной структуры]
Полный модуль новых стрелок допускает 10 комплексных интертвинеров в прямой
ориентации и 14 после вещественного завершения, поэтому общий коннектор неоднозначен.
Но на шестирёберной Ходж-опоре прямая часть равна нулю, а партнёр вещественной структуры
оставляет ровно один канал $Y_R\to Q_L\rightsquigarrow u_R$. Он уникален с
точностью до общего масштаба, физически типизирован и имеет ранг не выше
трёх.
\end{theorem}

\begin{nogocriterion}[Запрет повышения ранга оболочкой]
Нельзя объявлять одномерный триплетный канал полным коннектором
$M_{22}\leftrightarrow M_{21}$ только потому, что полная матричная оболочка
породил бы из любого ненулевого элемента весь угол. Следующий гейт должен
встроить именно физическое отображение ранга три в комплекс вещественной структуры и Ходжа и
проверить, какие центральные проекторы и блоки кривизны оно действительно
устраняет.
\end{nogocriterion}

\begin{protocol}
\begin{enumerate}
 \item полный модуль стрелок: \textbf{$36$ комплексных измерений};
 \item прямые физические интертвинеры: \textbf{$10$};
 \item после вещественного завершения: \textbf{$14$};
 \item прямые интертвинеры на выбранной опоре: \textbf{$0$};
 \item интертвинеры вещественной структуры на выбранной опоре: \textbf{$1$};
 \item единственное ребро: \textbf{$Q_L--Y_R$};
 \item компонента концевого сектора: \textbf{$u_R$};
 \item максимальный ранг: \textbf{$3$};
 \item полный коннектор носителей: \textbf{не получен};
 \item следующий гейт: \textbf{подъём вещественной структуры канала $Q_LY_R\to u_R$}.
\end{enumerate}
\end{protocol}

Машинный сертификат сохранён в
\path{s2t_v8_physical_arrow_endpoint_intertwiner_classification_gate_results.json}.

\begin{thebibliography}{9}
\bibitem{v8-intertwiner-krajewski}
T.~Krajewski, \emph{Classification of Finite Spectral Triples},
J. Geom. Phys. 28 (1998), 1--30; arXiv:hep-th/9701081.
\bibitem{v8-intertwiner-cacic}
B.~\v{C}a\'ci\'c, \emph{Moduli Spaces of Dirac Operators for Finite
Spectral Triples}, J. Noncommut. Geom. 3 (2009), 453--495;
arXiv:0902.2068.
\end{thebibliography}